In this section we give the basic concepts and definitions of coding theory (\cite{macwilliams1977theory}), required to formally define
the problems on which HQC and similar cryptosystems are based.

\subsection{General definitions}
We denote with $\mathds{F}_2$ the binary finite field and with $\mathds{Z}$ the integer ring.\\
We define $\mathcal{V} = \mathds{F}_2^n$ the vector space of dimension $n$ over $\mathds{F}_2$ 
and $\mathcal{R}=\mathds{F}_2[X]/(X^n-1)$ the ring of binary polynomials of degree less than $n$ for some $n \in \mathds{N}$.\\
Additionally, we denote by $\texttt{wt}(\cdot)$ the Hamming weight of a vector \textit{i.e.} the number of its non-zero coordinates.

\begin{definition}[Circulant matrix]
    Let $\mathbf{v} = (v_0, \cdots, v_{n-1}) \in \mathcal{V}$. The circulant matrix induced by $\mathbf{v}$ is defined as
    a matrix whose $i$-th row is the vector $\mathbf{v}$ shifted by $i$ positions to the right:
    \begin{equation*}
        \mathbf{circ(v)} = \begin{pmatrix}
            v_0 & v_1 & \cdots & v_{n-1} \\
            v_{n-1} & v_0 & \cdots & v_{n-2} \\
            \vdots & \vdots & \ddots & \vdots \\
            v_1 & v_2 & \cdots & v_0
            \end{pmatrix} \in \mathds{F}_2^{n\times n}
    \end{equation*}
\end{definition}
The algebra of $n\times n$ circulant matrices over $\mathds{F}_2$ is isomorphic to the algebra of polynomials
in the ring $\mathcal{R}$:
\begin{align*}
    \mathbf{circ(v)} \simeq v_0 + v_1x + v_2x^2 + \cdots + x_{n-1}x^{n-1} \\
\end{align*}

\begin{definition}[Linear code]
    A linear code $\mathcal{C}$ of lenght $n$ and dimension $k$ (denoted as $[n, k]$) is a subspace of $\mathcal{V}$ 
    of dimension $k$; the elements of $\mathcal{C}$ are called \textit{codewords}.
\end{definition}

\begin{definition}[Generator Matrix]
    Let $\mathcal{C} = [n,k]$ be a linear code. The matrix $\mathbf{G} \in \mathds{F}_2^{k\times n}$ is called the generator 
    matrix of $\mathcal{C}$ if
    \begin{equation*}
        \mathcal{C} = \{\mathbf{mG}, \text{ for } \mathbf{m}\in \mathds{F}_2^k\}
    \end{equation*}
\end{definition}

\begin{definition}[Parity-Check Matrix]
    Let $\mathcal{C} = [n,k]$ be a linear code. The matrix $\mathbf{H} \in \mathds{F}_2^{(n-k)\times n}$ is called the 
    parity-check matrix of $\mathcal{C}$ if
    \begin{equation*}
        \mathcal{C} = \{\mathbf{v} \in \mathds{F}_2^n \text{ such that } \mathbf{Hv}^\top = 0\}
    \end{equation*}
    In other words, $\mathbf{H}$ is the generator matrix of the dual code $\mathcal{C}^\perp$.
\end{definition}


\begin{definition}[Syndrome]
    Let $\mathcal{C} = [n,k]$ be a code, $\mathbf{H}$ its parity matrix and $\mathbf{v} \in \mathds{F}_2^n$ a word. 
    The syndrome of $\mathbf{v}$ is defined as
    \begin{equation*}
        \sigma(\mathbf{v}) = \mathbf{Hv}^\top
    \end{equation*}
    From this definition, we can say that a vector $\mathbf{v}$ is a codeword if and only if its syndrome is the zero vector.
\end{definition}

\begin{definition}[Minimum distance]
    Let $\mathcal{C} = [n,k]$ be a code and $\omega$ a norm on $\mathcal{V}$. The minimum distance of $\mathcal{C}$ is defined as
    \begin{equation*}
        d = \min_{\mathbf{v}, \mathbf{w} \in \mathcal{C}, \mathbf{v}\neq \mathbf{w}} \omega(\mathbf{v} - \mathbf{w})
    \end{equation*}
    A code with minimum distance $d$ ($\mathcal{C} = [n,k,d]$) can decode arbitrary codewords with at most
    $\delta = \lfloor \frac{d-1}{2}\rfloor$ errors.
\end{definition}

\begin{definition}[Quasi-cyclic Codes]
    Let $\mathbf{c} = (\mathbf{c}_0, \cdots, \mathbf{c}_{s-1}) \in \mathds{F}_2^{sn}$ be a vector, where we consider 
    each $\mathbf{c}_i$ as a vector in $\mathds{F}_2^n$.
    A linear code $\mathcal{C} = [sn,k,d]$ is Quasi-Cyclic of index $s$ if, for any $\mathbf{c} \in \mathcal{C}$, 
    the vector obtained by applying a circular shift to every block $\mathbf{c}_i$ simultaneously is also a codeword.
\end{definition}


\subsection{Difficult problems for cryptography}
The problems we are interested in are all variants of the \textit{decoding problem}, which requires to find 
the closest codeword to a given word with respect to a certain metric (the Hamming weight in our case).\\
In the case of linear codes we have shown that the decoding problem stays the same whether one is given 
the syndrome of the received vector or the vector itself; for this reason, we will only consider 
the \textit{Syndrome Decoding} (\textsf{SD}) problem.\\

\begin{definition}[\textsf{SD} Distribution]
    Given $n,k,w \in \mathds{N}$, the distribution $\textsf{SD}(n, k, w)$ randomly samples a parity-check matrix 
    $\mathbf{H} \xleftarrow{\mathdollar} \mathds{F}_2^{(n-k)\times n}$, a vector 
    $\mathbf{x} \xleftarrow{\mathdollar} \mathds{F}_2^n$ such that $\texttt{wt}(\mathbf{x}) = w$ 
    and returns ($\mathbf{H}, \sigma(\mathbf{x})$).
\end{definition}

\begin{definition}[\textsf{SD} Problem]
    Given $(\mathbf{H}, \mathbf{y}^\top) \in \mathds{F}_2^{(n-k)\times n} \times mathds{F}_2^{(n-k)}$ from the $\textsf{SD}(n, k, w)$ 
    distribution, the \textsf{SD} problem requires to find a vector $\mathbf{x} \in \mathds{F}_2^n$ such that 
    $\sigma(\mathbf{x}) = \mathbf{y}^\top$ and $\texttt{wt}(\mathbf{x}) = w$.
\end{definition}
This problem has been proven to be NP-complete in~\cite{berlekamp1978inherent}.

\begin{definition}[Decisional \textsf{SD} Problem]
    Given $(\mathbf{H}, \mathbf{y}^\top) \in \mathds{F}_2^{(n-k)\times n} \times mathds{F}_2^{(n-k)}$, the \textsf{DSD} problem 
    requires to determine whether $(\mathbf{H}, \mathbf{y}^\top)$ comes from the \textsf{SD} distribution or from the random
    distribution over $\mathds{F}_2^{(n-k)\times n} \times mathds{F}_2^{(n-k)}$.
\end{definition}


 



